\subsection{Modelo hidráulico y cálculo de la distribución}
\label{sec:man-asp-red}

Se identifica como $C$ la salida de alimentación central, situada a
$z_C=1{,}770$ m. Desde ella parten dos recorridos exteriores simétricos
y un recorrido interior. En cada lado exterior, $S1$ y $S2$ son las tomas
superiores; $E0$ es el inicio del tubo descendente y $E1$ a $E5$ son sus
tomas. En el ramal interior, $I0$ es la entrada y $I1$ a $I5$ son las
tomas ordenadas de arriba hacia abajo. Los puntos de la
Figura~\ref{fig:man-asp-esquema} son estaciones de cálculo, no un plano
de montaje.

\begin{figure}[htbp]
\centering
\begin{tikzpicture}[>=Stealth, every node/.style={font=\small},
 box/.style={draw,rounded corners=2pt,align=center,minimum height=10mm,text width=24mm},
 flow/.style={->,thick}]
\node[box] (tan) at (0,0) {Depósito\\abierto};
\node[box] (bom) at (3.3,0) {Bomba};
\node[box] (fil) at (6.6,0) {Filtro};
\node[box] (uni) at (9.9,0) {Unión\\rotativa};
\draw[flow] (tan)--node[above,font=\scriptsize]{DN50}(bom);
\draw[flow] (bom)--node[above,font=\scriptsize]{DN40}(fil);
\draw[flow] (fil)--node[above,font=\scriptsize]{DN40}(uni);
\node[box] (c) at (9.9,-1.8) {$C$\\Distribución};
\draw[flow] (uni)--node[right,font=\scriptsize]{DN40}(c);
\node[box,text width=34mm] (a) at (0,-3.6) {Lado exterior 1\\$S1,S2,E1\ldots E5$\\7 boquillas};
\node[box,text width=34mm] (b) at (4.9,-3.6) {Lado exterior 2\\Simétrico al lado 1\\7 boquillas};
\node[box,text width=34mm] (i) at (9.9,-3.6) {Ramal interior\\$I1\ldots I5$\\5 boquillas};
\draw[flow] (c.west)--(0,-1.8)--(a.north);
\draw[flow] (4.9,-1.8)--(b.north);
\draw[flow] (c)--(i);
\node[align=center,text width=12cm] at (4.9,-5) {Después de la aplicación, el líquido retorna al depósito.\\El dimensionamiento del retorno y del tanque se evalúa por separado.};
\end{tikzpicture}
\caption{Esquema funcional del circuito y de los grupos simultáneos. Elaboración propia; sin escala.}
\label{fig:man-asp-esquema}
\end{figure}

\subsubsection{Geometría hidráulica adoptada}
Se consideran diámetros exteriores de 33,4; 48,3 y 60,3 mm, con espesor
de 2,77 mm. El diámetro que interviene en la velocidad es el interior:
\begin{equation}
 D_{25}=33{,}4-2(2{,}77)=27{,}86\ \mathrm{mm},\quad
 D_{40}=42{,}76\ \mathrm{mm},\quad D_{50}=54{,}76\ \mathrm{mm}.
\end{equation}
Estas dimensiones se adoptan para el predimensionamiento y deben cotejarse
con la tubería suministrada. Se supone rugosidad $\varepsilon=0{,}045$ mm.

\subsubsection{Ecuaciones por tramo}
Para cada tubo se calculan la velocidad, Reynolds y factor de Darcy:
\begin{align}
 v&=\frac{4Q}{\pi D_i^2},\qquad Re=\frac{\rho vD_i}{\mu},
 \label{eq:man-asp-vre}\\
 f&=\frac{0{,}25}{\left[\log_{10}\left(
 \frac{\varepsilon}{3{,}7D_i}+\frac{5{,}74}{Re^{0{,}9}}\right)\right]^2}.
 \label{eq:man-asp-sj}
\end{align}
La Ecuación~\eqref{eq:man-asp-sj} es la aproximación de Swamee--Jain para
el régimen turbulento \cite{manSwamee}. Las pérdidas del tramo son
\begin{equation}
 h_L=h_f+h_m
 =\left(f\frac{L}{D_i}+\sum K_m\right)\frac{v^2}{2g}.
 \label{eq:man-asp-hl}
\end{equation}
Los coeficientes $K_m$ son adimensionales y se refieren a la velocidad del
tramo donde se colocan. No deben confundirse con el coeficiente de boquilla
$K_q$. Para el caso base se adoptan factores $f_T$ de 0,023; 0,020 y 0,019
para DN25, DN40 y DN50; se representan un paso de derivación por $20f_T$,
una salida de derivación por $60f_T$ y un codo por $30f_T$.
Son coeficientes de predimensionamiento, sujetos a confirmación con los
accesorios reales. Por ejemplo, la entrada a un lado exterior toma
$60(0{,}019)=1{,}14$, y la entrada interior,
$(60+30)(0{,}023)=2{,}07$.

\subsubsection{Relación entre presión, caudal y elevación}
En un tramo sin bomba, Bernoulli proporciona
\begin{equation}
 \frac{p_j}{\rho g}+z_j+\frac{v_j^2}{2g}
 =\frac{p_i}{\rho g}+z_i+\frac{v_i^2}{2g}-h_{L,ij}.
 \label{eq:man-asp-bernoulli-ramal}
\end{equation}
Para la primera aproximación de los colectores se desprecia la diferencia
de cargas cinéticas entre sus estaciones. Así,
\begin{equation}
 p_j\simeq p_i-\rho g\left[h_{L,ij}+(z_j-z_i)\right],\qquad
 q_j=K_q\sqrt{p_j},
 \label{eq:man-asp-recursion}
\end{equation}
con $p_j$ en bar únicamente en la segunda expresión, por las unidades
elegidas para $K_q$. La primera expresión emplea pascales.
La menor elevación incrementa la presión; la pérdida de carga la reduce.
Las tomas superiores se representan a 1,740 m, en el colector. El descenso
dentro del cuerpo y de las conexiones de cada boquilla no se modela.
Las diferencias cinéticas y estas conexiones se deben incluir al cerrar
la red; por ello se presenta un predimensionamiento, no una predicción
exacta de la instalación.

\subsubsection{Procedimiento de solución por aproximaciones sucesivas}
Se inicia con $q_j=9{,}11767$ L/min y $p_C=4$ bar. Se suman los caudales
aguas abajo de cada tramo, se calculan sus pérdidas y se obtienen las
presiones con la Ecuación~\eqref{eq:man-asp-recursion}. Después se actualiza
cada caudal con su presión y se ajusta la presión central mediante
\begin{equation}
 p_C^{(k+1)}=p_C^{(k)}+\left(4-p_{min}^{(k)}\right)
 \quad[\mathrm{bar}].
 \label{eq:man-asp-iteracion}
\end{equation}
Así se busca que la toma más desfavorable tenga 4 bar. Es una resolución
de ecuaciones de continuidad y energía; las sucesivas aproximaciones se
exponen para que el cálculo pueda comprobarse con una calculadora: la
Tabla~\ref{tab:man-asp-iteraciones} registra las aproximaciones y la
Tabla~\ref{tab:man-asp-tramos}, los tramos del estado convergido.

\begin{table}[htbp]
\centering\small
\caption{Aproximaciones sucesivas para determinar la presión central.}
\label{tab:man-asp-iteraciones}
\begin{tabular}{@{}rrrr@{}}
\toprule
Iteración & $p_C$ (bar) & $p_{\min}$ (bar) & $Q$ actualizado (L/min)\\\midrule
0 & 4{,}000000 & 4{,}001211 & 174{,}019197\\
1 & 3{,}998789 & 3{,}999984 & 173{,}991281\\
2 & 3{,}998805 & 4{,}000001 & 173{,}991697\\
3 & 3{,}998805 & 4{,}000000 & 173{,}991684\\
\bottomrule
\end{tabular}
\end{table}
En las dos últimas filas la diferencia de caudal total es menor que
0,001 L/min. Los dígitos adicionales permiten revisar el cierre aritmético;
no representan la precisión física de las hipótesis.

\begin{table}[htbp]
\centering\small
\caption{Datos por tramo y pérdidas del estado convergido.}
\label{tab:man-asp-tramos}
\begin{tabular}{@{}lrrrrrr@{}}
\toprule
Tramo & DN & $L$ (m) & $\sum K_m$ & $Q$ (L/min) & $h_L$ (m) & $z_f$ (m)\\\midrule
C--S1 & 50 & 0{,}5360 & 1{,}140 & 64{,}113 & 0{,}01436 & 1{,}7400\\
S1--S2 & 50 & 0{,}0880 & 0{,}380 & 54{,}995 & 0{,}00323 & 1{,}7400\\
S2--E0 & 50 & 0{,}2760 & 0{,}950 & 45{,}877 & 0{,}00577 & 1{,}7400\\
E0--E1 & 40 & 0{,}3940 & 0{,}400 & 45{,}877 & 0{,}00902 & 1{,}3460\\
E1--E2 & 40 & 0{,}0820 & 0{,}400 & 36{,}718 & 0{,}00415 & 1{,}2640\\
E2--E3 & 40 & 0{,}0820 & 0{,}400 & 27{,}550 & 0{,}00235 & 1{,}1820\\
E3--E4 & 40 & 0{,}0820 & 0{,}400 & 18{,}374 & 0{,}00105 & 1{,}1000\\
E4--E5 & 40 & 0{,}0300 & 0{,}400 & 9{,}188 & 0{,}00025 & 1{,}0700\\
C--I0 & 25 & 0{,}2800 & 2{,}070 & 45{,}766 & 0{,}18508 & 1{,}7400\\
I0--I1 & 25 & 0{,}4018 & 0{,}460 & 45{,}766 & 0{,}06526 & 1{,}3382\\
I1--I2 & 25 & 0{,}0954 & 0{,}460 & 36{,}630 & 0{,}02795 & 1{,}2428\\
I2--I3 & 25 & 0{,}0953 & 0{,}460 & 27{,}486 & 0{,}01581 & 1{,}1475\\
I3--I4 & 25 & 0{,}0953 & 0{,}460 & 18{,}334 & 0{,}00710 & 1{,}0522\\
I4--I5 & 25 & 0{,}0954 & 0{,}460 & 9{,}172 & 0{,}00181 & 0{,}9568\\
\bottomrule
\end{tabular}
\par\smallskip\footnotesize
Se muestra un lado exterior; el segundo es simétrico. $z_f$ es la elevación
final del tramo. Los caudales y pérdidas son resultados; $L$, DN y
$\sum K_m$ son entradas del modelo.
\end{table}

\subsubsection{Ejemplo completo: primer tramo del lado exterior}
El tramo $C$--$S1$ alimenta siete boquillas. En el estado convergido conduce
$Q=64{,}11274$ L/min $=1{,}06855\times10^{-3}$ m$^3$/s. Se tiene
$D_i=0{,}05476$ m, $L=0{,}536$ m y $\sum K_m=1{,}14$. Entonces:
\begin{align}
 v&=\frac{4(1{,}06855\times10^{-3})}{\pi(0{,}05476)^2}
 =0{,}45371\ \mathrm{m/s},\\
 Re&=\frac{982{,}120(0{,}45371)(0{,}05476)}{4{,}52269\times10^{-4}}
 =5{,}3952\times10^4,\\
 f&=\frac{0{,}25}{\left[\log_{10}\left(
 \frac{0{,}000045}{3{,}7(0{,}05476)}+
 \frac{5{,}74}{(53952)^{0{,}9}}\right)\right]^2}=0{,}023396,\\
 h_{L,C-S1}&=\left(0{,}023396\frac{0{,}536}{0{,}05476}+1{,}14\right)
 \frac{0{,}45371^2}{2(9{,}81)}=0{,}014363\ \mathrm{m}.
\end{align}
El segundo tramo, $S1$--$S2$, conduce el caudal restante después de la
primera boquilla y pierde 0,003230 m. Como $S2$ es la toma desfavorable,
la presión central se puede verificar directamente:
\begin{align}
 p_C&=p_{S2}+\rho g\left[h_{L,C-S1}+h_{L,S1-S2}+z_{S2}-z_C\right]\notag\\
 &=400000+982{,}120(9{,}81)
 \left(0{,}014363+0{,}003230+1{,}740-1{,}770\right)\notag\\
 &=\boxed{399880{,}47\ \mathrm{Pa}=3{,}99880\ \mathrm{bar}}.
 \label{eq:man-asp-pc}
\end{align}
El resultado es ligeramente inferior a 4 bar porque el descenso de 30 mm
aporta más presión hidrostática que la perdida en esos dos tramos. La
Tabla~\ref{tab:man-asp-boquillas} reúne la presión y el caudal de cada toma.

\begin{table}[htbp]
\centering\small
\caption{Presión y caudal calculados en las tomas de boquilla.}
\label{tab:man-asp-boquillas}
\begin{tabular}{@{}lrrrr@{}}
\toprule
Punto & Cantidad & $z$ (m) & $p_g$ (bar) & $q$ (L/min)\\\midrule
S1 & 2 & 1{,}7400 & 4{,}00031 & 9{,}11802\\
S2 & 2 & 1{,}7400 & 4{,}00000 & 9{,}11767\\
E1 & 2 & 1{,}3460 & 4{,}03653 & 9{,}15921\\
E2 & 2 & 1{,}2640 & 4{,}04404 & 9{,}16772\\
E3 & 2 & 1{,}1820 & 4{,}05171 & 9{,}17641\\
E4 & 2 & 1{,}1000 & 4{,}05951 & 9{,}18524\\
E5 & 2 & 1{,}0700 & 4{,}06238 & 9{,}18848\\
I1 & 1 & 1{,}3382 & 4{,}01629 & 9{,}13621\\
I2 & 1 & 1{,}2428 & 4{,}02279 & 9{,}14360\\
I3 & 1 & 1{,}1475 & 4{,}03044 & 9{,}15230\\
I4 & 1 & 1{,}0522 & 4{,}03894 & 9{,}16194\\
I5 & 1 & 0{,}9568 & 4{,}04796 & 9{,}17216\\
\bottomrule
\end{tabular}
\par\smallskip\footnotesize
La cantidad 2 corresponde a puntos equivalentes en los dos lados exteriores.
Los puntos $I1$--$I5$ pertenecen a un único ramal interior.
\end{table}

La suma de un lado exterior es 64,11274 L/min y la del ramal interior
45,76621 L/min. Por continuidad, el caudal final es
\begin{equation}
 Q=2(64{,}11274)+45{,}76621
 =\boxed{173{,}99168\ \mathrm{L/min}}
 =\boxed{2{,}89986\times10^{-3}\ \mathrm{m^3/s}}.
 \label{eq:man-asp-qfinal}
\end{equation}
La variación respecto de $Q_0=173{,}23564$ L/min es 0,44\,\%.
Las boquillas inferiores entregan algo más de líquido al recibir más
presión. La dispersión extrema de caudal predicha es
\begin{equation}
 \delta_q=100\left(\sqrt{\frac{4{,}06238}{4{,}00000}}-1\right)
 =0{,}777\ \mathord{\text{\%}}.
\end{equation}
Esta dispersión se refiere al caudal de las boquillas del modelo; no es
un indicador de uniformidad de limpieza sobre la pieza.

\subsection{Diámetro de alimentación y pérdidas de la línea principal}
\label{sec:man-asp-alimentacion}

Se adopta como criterio de predimensionamiento una velocidad de 1,5 a
3 m/s en la alimentación a presión. No se aplica ese intervalo como
restricción universal a todos los colectores. Para DN40,
\begin{equation}
 v_{40}=\frac{4(2{,}89986\times10^{-3})}{\pi(0{,}04276)^2}
 =\boxed{2{,}01935\ \mathrm{m/s}}.
 \label{eq:man-asp-v40}
\end{equation}
El Reynolds y el factor de Darcy son
\begin{align}
 Re_{40}&=\frac{982{,}120(2{,}01935)(0{,}04276)}{4{,}52269\times10^{-4}}
 =1{,}87507\times10^5,\\
 f_{40}&=\frac{0{,}25}{\left[\log_{10}\left(
 \frac{0{,}000045}{3{,}7(0{,}04276)}+
 \frac{5{,}74}{(187507)^{0{,}9}}\right)\right]^2}=0{,}021478.
\end{align}

\subsubsection{Línea estacionaria}
Se adopta una longitud de 14 m, seis codos, una válvula de retención y
una válvula de bola. La longitud y los accesorios son hipótesis de trazado.
Con $K_{codo}=0{,}60$, $K_{ret}=2{,}00$ y $K_{bola}=0{,}06$,
\begin{equation}
 \sum K_{est}=6(0{,}60)+2{,}00+0{,}06=5{,}66.
\end{equation}
La fricción distribuida y las pérdidas locales se calculan por separado:
\begin{align}
 h_{f,est}&=0{,}021478\frac{14}{0{,}04276}
 \frac{2{,}01935^2}{2(9{,}81)}=1{,}46154\ \mathrm{m},\\
 h_{m,est}&=5{,}66\frac{2{,}01935^2}{2(9{,}81)}=1{,}17636\ \mathrm{m},\\
 h_{L,est}&=1{,}46154+1{,}17636=\boxed{2{,}63790\ \mathrm{m}}.
 \label{eq:man-asp-hest}
\end{align}

\subsubsection{Alimentación del distribuidor}
Entre la unión rotativa y el punto $C$ se considera DN40, longitud de
1,75 m y un codo. Como el diámetro y el caudal son los mismos, se
mantienen $v_{40}$ y $f_{40}$:
\begin{align}
 h_{L,rot}&=\left(0{,}021478\frac{1{,}75}{0{,}04276}+0{,}60\right)
 \frac{2{,}01935^2}{2(9{,}81)}\notag\\
 &=0{,}18269+0{,}12470=\boxed{0{,}30740\ \mathrm{m}}.
 \label{eq:man-asp-hrot}
\end{align}

\subsubsection{Contraste con DN25}
Para conducir el mismo caudal por un diámetro interior de 27,86 mm,
\begin{equation}
 v_{25}=\frac{4(2{,}89986\times10^{-3})}{\pi(0{,}02786)^2}
 =4{,}75691\ \mathrm{m/s}.
\end{equation}
Se obtiene $Re=2{,}87789\times10^5$, $f=0{,}023022$ y
$\sum K=(6\times30+100+3)(0{,}023)=6{,}509$. Para los mismos 14 m,
\begin{equation}
 h_{L,25}=\left(0{,}023022\frac{14}{0{,}02786}+6{,}509\right)
 \frac{4{,}75691^2}{2(9{,}81)}=20{,}8499\ \mathrm{m}.
\end{equation}
DN40 reduce esa pérdida de 20,85 a 2,64 m y satisface el criterio de
velocidad adoptado. Por esta razón se conserva como diámetro de
predimensionamiento de la alimentación. La comprobación corresponde
a capacidad hidráulica, no al espesor resistente final del tubo.

\subsection{Línea de succión y condiciones del depósito}
\label{sec:man-asp-succion}

Se considera un depósito abierto, una línea de succión DN50 de 1,5 m
y un nivel de líquido 0,5 m por encima del eje de la bomba. Este nivel
es una condición que debe cumplir el diseño del tanque, no una
medición de planta. La velocidad y el factor de fricción son
\begin{align}
 v_1&=\frac{4(2{,}89986\times10^{-3})}{\pi(0{,}05476)^2}
 =1{,}23129\ \mathrm{m/s},\\
 Re_1&=\frac{982{,}120(1{,}23129)(0{,}05476)}{4{,}52269\times10^{-4}}
 =1{,}46417\times10^5,\qquad f_1=0{,}020945.
\end{align}
Se adoptan coeficientes de entrada 0,50, de codo $30(0{,}019)=0{,}57$
y de válvula de bola $3(0{,}019)=0{,}057$. Por ello,
\begin{equation}
 \sum K_s=0{,}50+0{,}57+0{,}057=1{,}127.
\end{equation}
La pérdida en succión resulta
\begin{align}
 h_{L,s}&=\left(0{,}020945\frac{1{,}5}{0{,}05476}+1{,}127\right)
 \frac{1{,}23129^2}{2(9{,}81)}\notag\\
 &=0{,}04433+0{,}08709=\boxed{0{,}13142\ \mathrm{m}}.
 \label{eq:man-asp-hs}
\end{align}
No se ha incluido una caída adicional por un filtro en succión. Si se
incorpora uno, su pérdida debe añadirse tanto al balance de energía como
a la evaluación de cavitación.

\subsection{Altura que debe suministrar la bomba}
\label{sec:man-asp-bomba}

La altura de bombeo expresa energía por unidad de peso del líquido. No
significa que la bomba esté instalada a esa altura sobre el piso.
Se obtiene aplicando Bernoulli entre la superficie libre del depósito,
denominada $S$, y la sección de alimentación $C$:
\begin{equation}
 \frac{p_{S,g}}{\rho g}+z_S+\frac{v_S^2}{2g}+H_b
 =\frac{p_{C,g}}{\rho g}+z_C+\frac{v_C^2}{2g}+\sum h_L.
 \label{eq:man-asp-bern-global}
\end{equation}
El depósito está abierto, de modo que $p_{S,g}=0$, y su sección se
considera suficientemente grande para tomar $v_S\simeq0$.
Se adopta el eje de bomba a $z_b=0{,}50$ m; con nivel inundado de 0,50 m,
la superficie queda a $z_S=1{,}00$ m. La sección $C$ está a 1,770 m
y se evalúa sobre la alimentación DN40, de velocidad $v_C=2{,}01935$ m/s.

Además de las pérdidas de tubería, se suponen caídas de 0,30 bar en la
unión rotativa y de 0,20 bar en el filtro. Estas dos pérdidas requieren
confirmación mediante curvas comerciales. Su altura equivalente es
\begin{equation}
 h_{JR+f}=
 \frac{(0{,}30+0{,}20)10^5}{982{,}120(9{,}81)}
 =5{,}18963\ \mathrm{m}.
 \label{eq:man-asp-accesorios}
\end{equation}
Despejando $H_b$ se obtiene la expresión de cálculo:
\begin{align}
 H_b={}&\frac{p_{C,g}}{\rho g}+(z_C-z_S)
 +h_{L,est}+h_{L,rot}+h_{L,s}+h_{JR+f}+\frac{v_C^2}{2g}.
 \label{eq:man-asp-hb}
\end{align}
La sustitución muestra de dónde proviene cada metro de altura:
\begin{align}
 H_b={}&\frac{399880{,}47}{982{,}120(9{,}81)}+(1{,}770-1{,}000)
 +2{,}63790+0{,}30740\notag\\
 &+0{,}13142+5{,}18963+\frac{2{,}01935^2}{2(9{,}81)}\notag\\
 ={}&41{,}50462+0{,}77000+2{,}63790+0{,}30740\notag\\
 &+0{,}13142+5{,}18963+0{,}20784\notag\\
 ={}&\boxed{50{,}74881\ \mathrm{m}}.
 \label{eq:man-asp-hbnum}
\end{align}
La Tabla~\ref{tab:man-asp-altura} ordena estas contribuciones.

\begin{table}[htbp]
\centering\small
\caption{Contribuciones a la altura de bombeo del caso base.}
\label{tab:man-asp-altura}
\begin{tabularx}{\linewidth}{@{}Xr@{}}
\toprule Contribución física & Altura (m)\\\midrule
Presión requerida en la alimentación central & 41,50462\\
Desnivel entre nivel libre y alimentación central & 0,77000\\
Fricción y accesorios de la línea estacionaria & 2,63790\\
Fricción y codo de la alimentación del distribuidor & 0,30740\\
Pérdidas de la succión & 0,13142\\
Unión rotativa y filtro & 5,18963\\
Carga cinética en la alimentación central & 0,20784\\\midrule
Total & 50,74881\\\bottomrule
\end{tabularx}
\end{table}

La presión requerida representa aproximadamente el 81,8\,\% de la
altura total. La demanda principal es, por tanto, proporcionar presión
al conjunto de boquillas. Las pérdidas, aunque menores, deben incluirse
para no subdimensionar la alimentación.
Las pérdidas en los ramales después de $C$ ya están incorporadas al
determinar $p_C$; no se vuelven a sumar en la
Ecuación~\eqref{eq:man-asp-hb}. Tampoco se añade otra vez la energía del
chorro, que procede de la presión aplicada a la boquilla.

La cota elevada de la unión rotativa, $z_{JR}=3{,}52$ m, no se suma como
un desnivel adicional permanente: el líquido vuelve a descender hasta
$C$. En el balance global se utiliza la diferencia entre extremos,
$z_C-z_S$. La elevación intermedia sí interviene al comprobar presiones
locales, purga de aire y llenado.

\subsubsection{Comprobación mediante las bridas de la bomba}
En la descarga, con brida a $z_b=0{,}50$ m y el mismo diámetro que la
alimentación, se requiere
\begin{align}
 p_{2,g}&=p_{C,g}+\rho g\left[h_{L,est}+h_{L,rot}+z_C-z_b\right]
 +\Delta p_{JR}+\Delta p_f\notag\\
 &=399880{,}47+982{,}120(9{,}81)
 (2{,}63790+0{,}30740+1{,}770-0{,}500)+50000\notag\\
 &=490493{,}17\ \mathrm{Pa}=4{,}90493\ \mathrm{bar}.
 \label{eq:man-asp-p2}
\end{align}
En la brida de succión, la presión manométrica resulta de descontar la
pérdida y la aceleración desde el depósito:
\begin{align}
 p_{1,g}&=\rho g\left(h_s-h_{L,s}-\frac{v_1^2}{2g}\right)\notag\\
 &=982{,}120(9{,}81)
 \left(0{,}50-0{,}13142-\frac{1{,}23129^2}{2(9{,}81)}\right)
 =2806{,}64\ \mathrm{Pa}.
\end{align}
Suponiendo las dos bridas a igual elevación,
\begin{equation}
 H_b=\frac{490493{,}17-2806{,}64}{982{,}120(9{,}81)}
 +\frac{2{,}01935^2-1{,}23129^2}{2(9{,}81)}
 =50{,}74881\ \mathrm{m}.
\end{equation}
La coincidencia de ambos balances comprueba que se han usado presiones
y cargas cinéticas con referencias consistentes.

\subsection{Potencia hidráulica, potencia al eje y potencia eléctrica}
\label{sec:man-asp-potencia}

La potencia hidráulica es la energía que recibe el líquido por unidad
de tiempo. Con el caudal y la altura calculados:
\begin{align}
 P_h&=\rho gQH_b\notag\\
 &=982{,}120(9{,}81)(2{,}89986\times10^{-3})(50{,}74881)
 =\boxed{1417{,}87\ \mathrm{W}}.
 \label{eq:man-asp-ph}
\end{align}
Se adoptan una eficiencia de bomba $\eta_p=0{,}55$ y una eficiencia
de motor $\eta_m=0{,}85$. Son hipótesis de predimensionamiento, que
deben reemplazarse por las del equipo seleccionado:
\begin{align}
 P_{eje}&=\frac{P_h}{\eta_p}
 =\frac{1417{,}87}{0{,}55}=\boxed{2577{,}95\ \mathrm{W}},\\
 P_{el}&=\frac{P_{eje}}{\eta_m}
 =\frac{2577{,}95}{0{,}85}=\boxed{3032{,}88\ \mathrm{W}}.
 \label{eq:man-asp-pel}
\end{align}
La potencia nominal de salida de un motor se contrasta con la demanda
al eje de la bomba, aproximadamente 2,58 kW, y con su intervalo real
de operación. Los 3,03 kW representan consumo eléctrico estimado,
no la potencia nominal de un motor ya seleccionado. Tampoco incluyen
calentamiento del baño, accionamiento del distribuidor, secado ni auxiliares.

Con precisión compatible con las hipótesis, el punto de predimensionamiento es
\begin{equation}
 \boxed{Q\approx10{,}4\ \mathrm{m^3/h},\qquad H_b\approx51\ \mathrm{m}}.
 \label{eq:man-asp-punto}
\end{equation}
Una bomba candidata debe evaluarse mediante su curva, eficiencia,
materiales y comportamiento de succión. El punto calculado no sustituye
esa selección.

\subsection{Comprobación de la disponibilidad de NPSH}
\label{sec:man-asp-npsh}

El NPSH disponible expresa la carga absoluta que queda sobre la presión
de vapor del líquido en la entrada de la bomba. Es una propiedad de la
instalación y debe compararse con el NPSH requerido por la bomba
\cite{white2016}. Para el escenario de emplazamiento a 2335 m se adopta
la aproximación troposférica de la atmósfera estándar \cite{white2016}:
\begin{align}
 p_{atm}&=101325\left(1-2{,}25577\times10^{-5}z_{sitio}\right)^{5{,}25588}\notag\\
 &=101325\left(1-2{,}25577\times10^{-5}(2335)\right)^{5{,}25588}
 =76243{,}85\ \mathrm{Pa}.
\end{align}
La altitud del sitio no debe confundirse con las cotas geométricas medidas
desde el piso de la cabina. Para depósito abierto,
\begin{align}
 NPSH_A&=\frac{p_{atm}-p_v}{\rho g}+h_s-h_{L,s}\notag\\
 &=\frac{76243{,}85-21867{,}00}{982{,}120(9{,}81)}
 +0{,}50-0{,}13142\notag\\
 &=5{,}64391+0{,}50000-0{,}13142
 =\boxed{6{,}01249\ \mathrm{m}}.
 \label{eq:man-asp-npsh}
\end{align}
Si el nivel libre descendiera hasta el eje de la bomba, manteniendo
el caudal y las pérdidas, $h_s=0$ y el NPSH disponible disminuiría a
5,51249 m. Se debe verificar la condición más desfavorable de nivel,
temperatura, caudal y ensuciamiento de los elementos de succión.
No se puede declarar un margen suficiente mientras no se conozcan
$NPSH_R$ y el criterio de margen aplicable al equipo seleccionado.